% Source: https://artofproblemsolving.com/community/c3h1437568_radius_of_a_sphere_inscribed_in_tetrahedron
% Stable identity: aops-topic-1437568-q01
Each face is an equilateral triangle of side $6$, with area $\dfrac{\sqrt3}{4}\cdot6^2=9\sqrt3$. The inscribed sphere touches every face, and the radius to the point of contact is perpendicular to that face, so the distance from $O$ to each face is $r$.

\begin{asy}
usepackage("amsmath");
size(5.5cm);
defaultpen(fontsize(10pt));
pair P(real x,real y,real z){return (x+0.35*y, z+0.25*y);}
real h=2*sqrt(6);
pair A=P(0,0,0),B=P(6,0,0),C=P(3,3*sqrt(3),0),D=P(3,sqrt(3),h),O=P(3,sqrt(3),h/4),G=P(3,sqrt(3),0);
draw(A--C,dashed); draw(B--C,dashed); draw(D--C,dashed);
draw(O--A,blue+dashed); draw(O--B,blue+dashed); draw(O--C,blue+dashed); draw(O--D,blue+dashed);
draw(O--G,red+linewidth(0.8));
draw(A--B--D--cycle,linewidth(0.9));
dot(A);dot(B);dot(C);dot(D);dot(O,red);dot(G,red);
label("$A$",A,SW); label("$B$",B,SE); label("$C$",C,2*E); label("$D$",D,N);
label("$O$",O,NE,red); label("$r$",G+(0,0.25),1.5*E,red); label("$6$",(A+B)/2,S);
label("\mbox{}",(min(currentpicture,user=true)+max(currentpicture,user=true))/2,black); // keeps the text after the figure black when LaTeX typesets the labels inline
\end{asy}

\textbf{(a)} Tetrahedron $OABC$ has base $ABC$ and height $r$, so its volume is $\dfrac{1}{3}\cdot9\sqrt3\cdot r=\boxed{3\sqrt3\,r}$.

\textbf{(b)} In the same way each of $OABD$, $OACD$, $OBCD$ has a face of area $9\sqrt3$ as base and height $r$: each has volume $\boxed{3\sqrt3\,r}$.

\textbf{(c)} The four small tetrahedra fill $ABCD$, so their volumes add to the volume of $ABCD$. The centroid of face $ABC$ is at distance $\dfrac{6}{\sqrt3}=2\sqrt3$ from each of $A$, $B$, $C$, so the height of $ABCD$ is $\sqrt{6^2-(2\sqrt3)^2}=\sqrt{24}=2\sqrt6$, and
\[
V_{ABCD}=\dfrac{1}{3}\cdot9\sqrt3\cdot2\sqrt6=6\sqrt{18}=18\sqrt2 .
\]
Hence $4\cdot3\sqrt3\,r=18\sqrt2$, and
\[
r=\dfrac{18\sqrt2}{12\sqrt3}=\dfrac{3\sqrt2}{2\sqrt3}=\boxed{\dfrac{\sqrt6}{2}} .
\]
(This is one quarter of the height $2\sqrt6$: $O$ is the centroid of the tetrahedron.)
